\chapter{基于溯源图的威胁检测综述}

\section{引言}

高级持续性威胁（advanced persistent threat, APT）通常不是单个恶意程序或单条异常日志，而是由初始入侵、权限维持、横向移动、凭据窃取、数据收集和外传等阶段组合而成的长期攻击过程。传统主机入侵检测常依赖签名、单点异常或短窗口序列特征，能够发现已知恶意样本，却难以回答三个更重要的问题：攻击从哪里开始、经过了哪些实体、造成了哪些影响。系统溯源图（system provenance graph）为这一问题提供了结构化表达：进程、文件、套接字、注册表项等系统实体被建模为节点，读、写、执行、连接等事件被建模为带时间和类型的有向边，从而将审计日志转换为可追踪的信息流和控制流。

基于溯源图的入侵检测系统（provenance-based intrusion detection systems, PIDS）因此成为APT检测的重要方向。综述性研究通常将该方向分解为数据采集、图构建与缩减、图存储与查询、威胁检测、威胁狩猎和取证分析等环节\citep{zipperle2023survey,leng2022review}。随着机器学习和图学习方法进入该领域，研究重点又从规则匹配和统计异常转向节点级检测、自监督表示学习、在线推理和可解释归因\citep{li2025intelligent}。近两年，大语言模型（large language models, LLMs）进一步把外部威胁情报、自然语言攻击报告和溯源图行为解释纳入检测流程，形成了LLM+RAG、图语言对比学习和智能取证代理等新分支\citep{omniseq2025,provseek2025,aptcglp2025}。

本文围绕“基于溯源图的威胁检测”展开综述。与只罗列方法不同，本文希望建立一条可理解的演进线索：早期方法强调“能否长期记录并发现异常”，中期方法强调“能否在大图上准确定位恶意节点”，近期方法强调“能否在线运行、降低人工分析成本，并解释攻击语义”。为此，本文以本地选定的14篇论文为核心阅读材料，辅以近期公开论文清单和原始论文页面补充前沿趋势\citep{xythink2026papers}，从检测范式、评测体系和未解决问题三个层面进行梳理。

\section{溯源图检测任务与系统框架}

溯源图检测任务的输入通常来自Linux Audit、ETW、DTrace、CamFlow或EDR系统采集的审计事件。事件经过解析、去重、实体归并和因果方向恢复后，形成随时间增长的动态图。一个典型PIDS流水线包括四步：首先采集系统事件并进行规范化；其次通过缩减、压缩或窗口划分控制图规模；然后对节点、边、路径、子图或整图进行异常评分；最后将告警转化为安全分析员可以验证的攻击路径、根因节点和影响范围。

不同方法的关键差异在于检测粒度和先验假设。规则型系统依赖攻击模式、威胁情报或领域查询语言，优点是可解释、实时性强，缺点是难以发现未知攻击。异常型系统学习正常行为分布，适合零日和变体攻击，但容易把环境变化误判为攻击。学习型系统进一步利用GNN、Transformer、自监督编码器或强化学习学习图结构和时间语义，在精度上取得显著提升，却带来训练成本、泛化能力和解释性问题。近期LLM方法试图在告警之后引入外部知识库、威胁情报和自然语言解释，以降低人工调查成本。

\begin{figure}[htbp]
    \centering
    \begin{tikzpicture}[
        node distance=1.15cm,
        stage/.style={rounded corners, draw, align=center, text width=3.0cm, minimum height=1.05cm, fill=blue!8},
        trend/.style={rounded corners, draw, align=center, text width=3.0cm, minimum height=1.05cm, fill=green!8},
        arrow/.style={-Latex, thick}
    ]
        \node[stage] (rule) {规则/图匹配\\SLEUTH, HOLMES};
        \node[stage, right=of rule] (stat) {统计异常/图草图\\Unicorn};
        \node[stage, right=of stat] (gnn) {节点级GNN\\THREATRACE};
        \node[stage, below=of gnn] (self) {自监督与表示学习\\MAGIC, FLASH};
        \node[stage, left=of self] (online) {在线细粒度检测\\KAIROS, NODLINK};
        \node[trend, left=of online] (llm) {语义增强与LLM\\OMNISEC, OCR-APT};
        \node[trend, below=of online] (future) {新趋势\\RL, EDR, 数据集更新};
        \draw[arrow] (rule) -- (stat);
        \draw[arrow] (stat) -- (gnn);
        \draw[arrow] (gnn) -- (self);
        \draw[arrow] (self) -- (online);
        \draw[arrow] (online) -- (llm);
        \draw[arrow] (llm) -- (future);
    \end{tikzpicture}
    \caption{基于溯源图的威胁检测方法演进脉络}
    \label{fig:evolution}
\end{figure}

\section{代表性方法演进}

\subsection{从规则匹配到统计异常}

早期PIDS的直接目标是利用溯源图重建攻击因果链。基于规则或查询的方法把攻击知识编码为标签传播、图模式或领域语言查询，能够较快定位已知战术技术过程（tactics, techniques and procedures, TTPs）。中文综述将这类方法归纳为基于规则、基于异常和基于学习的威胁检测，并指出规则类方法适合威胁狩猎和取证，但依赖专家知识\citep{leng2022review}。Zipperle等则强调，PIDS除了检测算法外，还必须面对采集开销、图缩减、图存储、查询和benchmark缺失等系统问题\citep{zipperle2023survey}。

Unicorn是统计异常阶段的重要代表。它面向APT“低速慢行”的特点，用图草图和固定大小的长期结构摘要表示系统行为，在不保存完整图的情况下支持运行时异常检测\citep{unicorn2020}。这类方法的价值在于证明长时间上下文确实有助于区分隐蔽攻击与正常行为；局限在于检测粒度偏粗，通常只能给出图级或粗粒度告警，后续攻击归因仍然依赖人工追踪。

\subsection{节点级图学习与归因质量}

随着GNN进入PIDS，研究开始从“整图异常”转向“节点级恶意实体定位”。THREATRACE把GraphSAGE用于溯源图节点表示学习，以无先验模式的方式识别异常实体，并强调实时监控和早期检测\citep{threatrace2022}。它解决了图级方法对少量恶意节点不敏感的问题，但仍主要依赖局部结构和系统调用分布，面对语义相近但安全含义不同的行为时可能产生误报。

KAIROS进一步从系统设计维度审视PIDS，提出实践系统应同时满足跨应用边界检测、攻击无关性、及时性和高效调查四个维度，并使用时序因果GNN进行检测与调查\citep{kairos2024}。ORTHRUS则把“归因质量”（quality of attribution, QoA）作为核心指标，指出接近完美的检测率并不等于可部署；如果告警包含大量与攻击无关的节点，安全分析员仍会陷入告警疲劳\citep{orthrus2025}。这一转向很关键：PIDS的评价对象不应只是模型输出的准确率，还应包括人工调查所需检查的节点数、攻击路径紧凑性和根因解释质量。

\subsection{自监督表示学习与在线检测}

MAGIC代表了自监督图表示学习路线。它通过掩码图表示学习缓解攻击样本稀缺问题，在只依赖正常数据或有限标签的情况下学习溯源图语义\citep{magic2024}。FLASH则指出现有方法常忽略命令行、文件路径、进程名、IP地址等语义属性，以及事件的时间因果顺序，因此结合Word2Vec语义编码和GNN上下文编码，并用嵌入缓存降低运行时推理开销\citep{flash2024}。这说明图结构本身不是全部，系统实体名称、路径和命令参数往往承载关键安全语义。

在线检测方向关注另一个实际问题：审计日志持续流入，溯源图规模巨大，如果每次都在全图上推理就难以部署。NODLINK将在线细粒度检测建模为在线Steiner Tree问题，利用内存缓存和近似子图搜索减少需要分析的区域，从而在开放世界场景下定位APT活动\citep{nodlink2024}。这类方法的共同趋势是用结构缩减、窗口划分、缓存或近似搜索在准确率和延迟之间折中。

\subsection{LLM辅助、强化学习与工程化落地}

LLM方法并不是简单替换检测模型，而是在三个位置补足传统PIDS短板。第一，OMNISEC使用无损缩减去除冗余信息流，构建良性行为知识库和威胁情报知识库，再用RAG提示LLM判断异常节点是否是真实攻击，并重建攻击图\citep{omniseq2025}。第二，OCR-APT结合子图异常检测和LLM，把低层审计日志转化为较完整的APT故事，提升取证叙事能力\citep{ocrapt2025}。第三，ProvSEEK和SHIELD等外部前沿工作将RAG、CoT和多智能体机制引入溯源取证，目标是生成可验证的调查摘要并降低误报解释成本\citep{provseek2025,shield2025}。

强化学习和工程化是另一条新线。Slot将检测过程建模为图强化学习，让策略在动态溯源图中学习自适应检测动作\citep{slot2025}。SYSARMOR则把溯源分析纳入EDR实践，通过微服务架构、实时流处理和异步检测引擎组合Falco规则、NODLINK和KNOWHOW，展示了PIDS从论文原型走向企业部署的路径\citep{sysarmor2026}。Auto-Prov进一步把LLM用于异构日志到功能嵌入溯源图的自动构建，说明未来改进点不只在检测模型，也在图构建和语义补全阶段\citep{autoprov2026}。

\section{核心论文精读与比较}

表\ref{tab:papers}总结了本地14篇论文及若干前沿补充文献的角色。可以看到，方法演进不是单纯从浅层模型到深层模型，而是围绕三个矛盾反复推进：检测未知攻击与控制误报之间的矛盾，大规模图分析与实时性之间的矛盾，高性能指标与可调查告警之间的矛盾。

\begin{table}[htbp]
    \centering
    \scriptsize
    \caption{代表性论文的一句话精读总结}
    \label{tab:papers}
    \renewcommand{\arraystretch}{1.05}
    \begin{tabularx}{\textwidth}{p{1.85cm} p{1.75cm} X X}
        \toprule
        论文 & 类型 & 核心动机与方法 & 主要启示 \\
        \midrule
        Unicorn & 统计异常 & 用图草图摘要长期系统行为，实现运行时APT异常检测。 & 长期上下文有价值，但图级告警归因能力有限。 \\
        THREATRACE & 节点级GNN & 用GraphSAGE学习实体角色，检测少量威胁节点。 & 检测粒度从图级推进到节点级。 \\
        KAIROS & 实用PIDS & 用时序因果GNN兼顾检测和调查。 & 实用系统需同时考虑范围、未知攻击、及时性和调查效率。 \\
        NODLINK & 在线检测 & 将细粒度APT检测转化为在线Steiner Tree近似搜索。 & 在线场景需要把全图问题压缩为局部可处理子图。 \\
        MAGIC & 自监督学习 & 用掩码图表示学习缓解攻击样本不足。 & 自监督预训练是减少标签依赖的重要方向。 \\
        FLASH & 语义+GNN & 融合Word2Vec语义编码、GNN结构编码和嵌入缓存。 & 文件路径、命令行和时间顺序是降低误报的关键语义。 \\
        ORTHRUS & 高质量归因 & 提出QoA视角，减少人工需要检查的无关节点。 & 告警质量应和准确率一样成为核心指标。 \\
        OMNISEC & LLM+RAG & 用良性行为库和威胁情报库辅助LLM判断异常行为。 & LLM适合作为语义判别和攻击链重建层。 \\
        OCR-APT & 子图异常+LLM & 先发现异常子图，再用LLM重建APT故事。 & 取证叙事是检测结果走向分析员的桥梁。 \\
        Slot & 强化学习 & 用图强化学习进行自适应APT检测。 & 动态策略可用于应对攻击过程和环境变化。 \\
        Bilot等 & SoK/复现 & 统一实现八个PIDS并指出九类实践短板。 & 复杂模型未必优于简单基线，评测方法需要重审。 \\
        Zipperle等 & 综述 & 从采集、缩减、检测和数据集系统化综述PIDS。 & PIDS是系统工程，不能只看检测算法。 \\
        冷涛等 & 中文综述 & 建立威胁发现、威胁狩猎和取证分析框架。 & 中文语境下清晰区分检测、狩猎和取证任务。 \\
        李振源等 & 智能溯源综述 & 比较机器学习方法的准确率、效率、鲁棒性和解释性。 & AI带来新能力，也引入对抗、投毒和部署成本问题。 \\
        \bottomrule
    \end{tabularx}
\end{table}

从表中可以进一步抽象出三种技术路线。第一类是结构驱动路线，代表为Unicorn、NODLINK和ProHunter，关注如何在大规模动态溯源图中保留关键因果关系并快速搜索可疑子图\citep{prohunter2026}。第二类是表示学习路线，代表为THREATRACE、KAIROS、MAGIC、FLASH和ORTHRUS，关注如何把节点属性、邻域结构和时间顺序编码成可学习表示。第三类是语义增强路线，代表为OMNISEC、OCR-APT、ProvSEEK、SHIELD、Auto-Prov和APT-CGLP，关注如何把威胁情报、自然语言报告和LLM推理接入检测与调查流程\citep{aptcglp2025,autoprov2026}。

\section{评测指标与benchmark}

PIDS评测常见指标包括Precision、Recall、F1、AUC、FPR和MCC。对于节点级方法，还需要统计TP、FP、TN、FN对应的节点数量；对于路径或子图级方法，需要评估攻击覆盖率、无关节点比例和路径完整性；对于在线系统，还需要报告吞吐、延迟、内存、CPU和训练/测试时间。Bilot等指出，很多论文报告的高性能来自不一致的评测设定，例如阈值选择、测试数据泄露、标签粒度不同或将攻击后代节点全部视为恶意，从而夸大了实际检测能力\citep{bilot2025simpler}。

\begin{table}[htbp]
    \centering
    \small
    \caption{常用数据集与适用场景}
    \label{tab:datasets}
    \renewcommand{\arraystretch}{1.18}
    \begin{tabularx}{\textwidth}{p{2.6cm} p{2.4cm} X X}
        \toprule
        数据集 & 主要系统/平台 & 常见用途 & 局限 \\
        \midrule
        DARPA TC E3/E5 & Linux, FreeBSD, Windows等 & PIDS最常用benchmark，包含Trace、Theia、Cadets、FiveDirections等子集。 & 攻击脚本化、时间跨度有限，标签和文档粒度不足。 \\
        DARPA OpTC & Windows 10企业环境 & 用于评估大规模Windows审计日志和多主机场景。 & 数据极大，处理门槛高，攻击占比极低。 \\
        StreamSpot & 信息流图 & 常用于流式图异常检测和图级基线比较。 & 场景较简单，与现代APT差距较大。 \\
        Unicorn数据集 & 企业级端点模拟 & 用于评估长期APT和供应链式场景。 & 规模和多样性仍不足以覆盖真实企业环境。 \\
        下一代数据集工作\citep{nextgendatasets2026} & 自动化攻击仿真 & 面向新型APT技术、细粒度标签和更完整攻击链。 & 仍处于建设阶段，需要社区共同验证。 \\
        \bottomrule
    \end{tabularx}
\end{table}

从benchmark角度看，DARPA TC仍是主流，但它也逐渐成为瓶颈。许多系统都在同一批数据上迭代，容易出现对特定数据集过拟合的现象。OpTC提供了更接近企业规模的Windows日志，但数据量和标注复杂度使复现实验门槛显著提高。NextGenDatasets等2026年工作明确指出，现有数据集存在攻击技术过时、时间尺度短、攻击链碎片化和标签不清晰等问题，并尝试通过自动化攻击仿真构建更大规模、更细粒度标注的新数据集\citep{nextgendatasets2026}。因此，未来PIDS的评测应同时包含旧benchmark上的可比性和新场景上的泛化性。

\section{存在问题与研究展望}

\subsection{误报、归因质量与人机协同}

异常检测的核心难题是正常行为本身会变化。软件更新、用户习惯变化、业务迁移和系统配置调整都会改变溯源图分布，导致模型把概念漂移误判为攻击。METANOIA提出以终身学习缓解概念漂移，并通过伪边、可疑状态转移和路径级过滤减少误报\citep{metanoia2025}。这表明未来检测系统不应是一次训练后长期固定的模型，而应能在不学习恶意行为的前提下持续吸收新正常行为。

同时，检测正确不等于调查可用。ORTHRUS和Bilot等工作的共同提醒是，研究需要从“是否命中攻击”扩展到“分析员需要付出多少努力”。一个可行方向是构造多层告警输出：第一层给出少量高置信异常节点，第二层给出从入口到影响点的紧凑攻击路径，第三层提供自然语言解释和证据链链接。LLM可以参与第三层解释，但必须由可验证的溯源查询结果约束，避免生成无法回溯的幻觉结论。

\subsection{跨系统泛化与语义增强}

不同操作系统、日志采集器和企业软件栈产生的实体类型、事件字段和命名习惯差异很大。仅依赖结构特征的GNN容易学习到数据集特有模式，而不是攻击语义。FLASH将文件路径、命令行和IP等语义属性纳入编码，APT-CGLP则尝试通过图语言对比预训练弥合CTI报告和溯源图之间的模态差距\citep{flash2024,aptcglp2025}。未来更可行的路线是“结构不变性+语义补全”：图结构负责表达因果关系，语义编码负责解释实体功能，跨域训练负责约束模型不要过度依赖单一平台特征。

Auto-Prov把这一问题前移到图构建阶段：它用LLM从异构日志中自动抽取溯源边和实体功能属性，再把功能上下文嵌入图中\citep{autoprov2026}。这提示我们，PIDS未来的竞争点可能不只在检测器，而在端到端数据管线：从日志解析、实体规范化、功能语义补全、图缩减到检测解释，每一步都会影响最终质量。

\subsection{面向未来的可能方案}

结合上述脉络，本文认为一个值得探索的方案是“语义增强的持续学习PIDS”。其基本思路如下：首先，在数据层保留可验证的溯源因果边，同时对进程名、文件路径、命令行参数和网络端点进行功能语义编码；其次，在模型层采用自监督图表示学习获得基础编码，再引入增量学习机制适应概念漂移；再次，在告警层使用路径级归因指标约束输出规模，只保留与攻击因果链强相关的节点和边；最后，在解释层使用RAG式LLM读取威胁情报和本次告警证据，生成带引用证据的调查摘要。

这一方案的关键不是让LLM直接判断“是否攻击”，而是让LLM在受约束的证据空间中完成语义补全和报告生成。检测决策仍应由可复现实验验证的图模型、统计模型或规则模型承担。这样既能利用LLM的语义理解能力，又能避免把安全决策完全交给不稳定的黑箱推理。评价时应同时报告检测指标、归因质量、人工检查节点数、在线延迟和跨数据集泛化性能。

\section{结论}

基于溯源图的威胁检测已经从规则匹配和统计异常发展到图神经网络、自监督学习、在线细粒度检测、LLM辅助调查和EDR工程化落地。总体趋势可以概括为三点：第一，检测粒度从图级异常走向节点、路径和子图级归因；第二，系统目标从追求单一准确率走向同时优化实时性、误报、解释性和人工工作量；第三，研究边界从检测模型扩展到数据集建设、日志到图的自动构建、语义补全和企业部署。

然而，该领域距离成熟落地仍有距离。现有benchmark过度集中于DARPA系列，攻击技术和业务场景更新不足；许多方法在统一评测框架下表现并不稳定；机器学习方法面临概念漂移、数据投毒和可解释性压力；LLM方法虽然改善语义解释，却必须解决事实约束、成本和隐私问题。未来研究应围绕可复现评测、跨系统泛化、高质量归因和人机协同调查展开，使溯源图真正成为连接底层审计日志和高层安全决策的可靠桥梁。
